\subsection{UNIVERSAL PROBLEM FOR DERIVATIONS; NON-COMMUTATIVE CASE}

Throughout the rest of §10 all the algebras are assumed to be associative and unital and all the algebra homomorphisms are assumed to be unital.

Let A be a K-algebra; the tensor product \(A \otimes_{K} A\) has canonically an (A, A)-bimodule structure under which

\[
x. (u \otimes v). y = (x u) \otimes (v y)\tag{40}
\]

for all \(x, y, u, v\) in A (§ 4, no. 3, Example 2). The K-linear mapping \(m: \mathrm{A} \otimes_{\mathbb{K}} \mathrm{A} \to \mathrm{A}\) corresponding to multiplication in A (and hence such that \(m(x \otimes y) = xy\) ) is an (A, A)-bimodule homomorphism; its kernel I is therefore a sub-bimodule of \(\mathrm{A} \otimes_{\mathbb{K}} \mathrm{A}\) .

Lemma 1. The mapping \(\delta_{\mathbf{A}}: x \mapsto x \otimes 1 - 1 \otimes x\) is a K-derivation of A into I and I is generated, as a left A-module, by the image of \(\delta_{\mathbf{A}}\) .

The first assertion follows from the fact that

\[
(x y) \otimes 1 - 1 \otimes (x y) = (x \otimes 1 - 1 \otimes x). y + x. (y \otimes 1 - 1 \otimes y)
\]

by (40). On the other hand, if the element \(\sum_{i} x_{i} \otimes y_{i}\) (for \(x_{i}, y_{i}\) in A) belongs to I, by definition \(\sum_{i} x_{i} y_{i} = 0\) and hence

\[
\sum_ {i} \left(x _ {i} \otimes y _ {i}\right) = \sum_ {i} x _ {i} (1 \otimes y _ {i} - y _ {i} \otimes 1)
\]

by (40), which completes the proof of the lemma.

PROPOSITION 17. The derivation \(\delta_{\mathbf{A}}\) has the following universal property: for every (A, A)-bimodule E and every K-derivation \(d: \mathbf{A} \to \mathbf{E}\) , there exists one and only one (A, A)-bimodule homomorphism \(f: \mathbf{I} \to \mathbf{E}\) such that \(d = f \circ \delta_{\mathbf{A}}\) .

Note first that, for every (A, A)-bimodule homomorphism \(f: I \to E, f \circ \delta_{A}\) is a derivation (no. 7, Proposition 9). Conversely, let \(d: A \to E\) be a K-derivation; then we prove first that if there exists an (A, A)-bimodule homomorphism \(f: I \to E\) such that \(d = f \circ \delta_{A}, f\) is uniquely determined by this condition for the definition of \(\delta_{A}\) gives

\[
f (x \otimes 1 - 1 \otimes x) = d x
\]

and our assertion follows from the fact that the image of \(\delta_{\mathbf{A}}\) already generates I as a left A-module (Lemma 1): hence of necessity

\[
f \left(\sum_ {i} x _ {i} \otimes y _ {i}\right) = \sum_ {i} x _ {i}. f (1 \otimes y _ {i} - y _ {i} \otimes 1) = - \sum_ {i} x _ {i}. d y _ {i}.
\]

Conversely, as the mapping \((x,y)\mapsto -x.dy\) of \(\mathbf{A}\times \mathbf{A}\) into \(\mathbf{E}\) is K-bilinear, there exists one and only one K-linear mapping \(g:\mathbf{A}\otimes_{\mathbb{K}}\mathbf{A}\rightarrow \mathbf{E}\) such that \(g(x\otimes y) = -x.dy\) ; it suffices to verify that the restriction \(f\) of \(g\) to \(\mathbf{I}\) is A-linear for the left and right A-module structures. The first assertion is obvious since \((xx') .dy = x.(x'.dy)\) ; to prove the second, note that, if \(\sum_{i}x_{i}\otimes y_{i}\in \mathbf{I}\) and \(x\in \mathbf{A}\) , then

\[
\sum_ {i} x _ {i}. d (y _ {i} x) = \sum_ {i} x _ {i}. d y _ {i}. x + \sum_ {i} (x _ {i} y _ {i}). d x
\]

but since \(\sum_{i}x_{i}y_{i} = 0\) by definition of I, this completes the proof.

We have thus defined a canonical K-module isomorphism \(f \mapsto f \circ \delta_{A}\)

\[
\operatorname{Hom} _ {(\mathbf {A}, \mathbf {A})} (\mathbf {I}, \mathbf {E}) \rightarrow \mathbf {D} _ {\mathbf {K}} (\mathbf {A}, \mathbf {E})
\]

the left hand side being the K-module of (A, A)-bimodule homomorphisms of A into E.
% semantic-fragments: legacy-input-seam
\relax{}
